% ─────────────────────────────────────────────────────────────────────
% AIME  — Appendix / Supporting Information (appendix.tex)
% Last updated: April 3rd, 2026
% ─────────────────────────────────────────────────────────────────────

\documentclass[11pt]{aime}
\usepackage{lmodern}

\title{Supporting Information for:\\A unified experimental-simulation dataset and surrogate models for surfactant property prediction}

\author[1]{Richard Beckmann}
\author[1]{Pablo Martínez Crespo}
\author[2]{Robert S. Jordan}
\author[3]{Marisa Gliege}
\author[4]{Santiago Miret}
\author[5]{Vijay Kris Narasimhan}
\author[1]{Rocío Mercado}

\affiliation[1]{AI Laboratory for Molecular Engineering (AIME), Department of Computer Science and Engineering, Chalmers University of Technology \& University of Gothenburg, Gothenburg, Sweden}
\affiliation[2]{Technology Research, Intel Corporation}
\affiliation[3]{Chief Technology Office, EMD Electronics}
\affiliation[4]{Lila Sciences}
\affiliation[5]{Merck KGaA, Darmstadt, Germany}

\contribution[*]{Corresponding author}

\correspondence{\email{rocio.mercado@chalmers.se}}

\abstract{%
This document contains supporting information for the main manuscript.
Section~\ref{app:shap} provides a detailed SHAP-based feature attribution analysis across all nine target properties, characterising which RDKit molecular descriptors drive model predictions, the degree of cross-target overlap among the most important features, and how quickly SHAP importance decays as a function of feature rank.
Section~\ref{app:ablation} reports the feature representation ablation study, comparing raw RDKit descriptors against Morgan circular fingerprints and PCA-compressed variants.
Section~\ref{app:feature_importance} presents the permutation feature importance analysis and the results of structure--property relationship analysis across the nine targets.
Section~\ref{app:unc_diagnostics} provides diagnostic scatter plots for the nearest-neighbour distance and ensemble disagreement uncertainty heuristics.
Section~\ref{app:conformal} gives a detailed analysis of the conformal calibration curves, including a discussion of the three calibration regimes and their practical implications for users.
}

\begin{document}
\maketitle

% =====================================================================
% SHAP FEATURE ATTRIBUTION
% =====================================================================

\section{SHAP-based feature attribution}
\label{app:shap}

To characterise which molecular descriptors drive model predictions, we computed SHAP (SHapley Additive exPlanations) values for all nine target models using a TreeExplainer applied to the XGBoost ensembles.
For each target, SHAP values were computed on the test set molecules across all five test splits and all five fold models within each split;
the mean absolute SHAP value per feature was then averaged across splits.

\begin{figure}[ht]
    \centering
    \includegraphics[width=\textwidth]{../surrogate-models/shap_per_target_ranked.png}
    \caption{Top-15 features ranked by mean $|\text{SHAP}|$ for each of the nine target properties.
    Bar colours indicate cross-target sharing: dark blue --- features appearing in the top-20 of three or more targets;
    light blue --- two targets;
    grey --- target-specific.
    The RDKit descriptor names follow the standard RDKit nomenclature.}
    \label{fig:shap_per_target}
\end{figure}

\begin{figure}[ht]
    \centering
    \includegraphics[width=\textwidth]{../surrogate-models/shap_presence_matrix.png}
    \caption{Binary presence matrix showing which RDKit features appear in the top-20 SHAP importance list for each target.
    Each row is one feature;
    each column is one target.
    Filled cells indicate membership in the top-20 list;
    the right-hand column reports the count of targets in which the feature appears.}
    \label{fig:shap_presence}
\end{figure}

\begin{figure}[ht]
    \centering
    \includegraphics[width=0.75\textwidth]{../surrogate-models/shap_feature_count_dist.png}
    \caption{Distribution of cross-target sharing among the top-20 SHAP features.
    Each bar shows the number of distinct RDKit descriptors that appear in the top-20 list for exactly $k$ targets ($k = 1, \ldots, 9$).
    Most features are target-specific;
    a small subset recurs across three or more models.}
    \label{fig:shap_count_dist}
\end{figure}

\begin{figure}[ht]
    \centering
    \includegraphics[width=\textwidth]{../surrogate-models/shap_crossover_heatmap.png}
    \caption{Heatmap of normalised mean $|\text{SHAP}|$ for features that appear in the top-20 list of at least two targets.
    Values within each target column are normalised so that the maximum is 1.
    The label suffix $\times k$ indicates the number of targets for which the feature is in the top-20.
    Features are sorted by the number of targets in which they appear.}
    \label{fig:shap_heatmap}
\end{figure}

\begin{figure}[ht]
    \centering
    \includegraphics[width=0.75\textwidth]{../surrogate-models/shap_universal_features.png}
    \caption{Grouped bar chart of features appearing in the top-20 SHAP list for three or more targets (``universal'' features).
    Each group of bars shows the normalised mean $|\text{SHAP}|$ for one feature across the targets in which it appears.}
    \label{fig:shap_universal}
\end{figure}

\begin{figure}[ht]
    \centering
    \includegraphics[width=\textwidth]{../surrogate-models/shap_decay_grid.png}
    \caption{Decay of mean $|\text{SHAP}|$ as a function of feature rank (log scale) for each of the nine targets.
    The shaded region marks the top-20 features.
    Dashed vertical lines indicate the ranks at which 50\% and 80\% of the cumulative SHAP importance is captured.}
    \label{fig:shap_decay_grid}
\end{figure}

\begin{figure}[ht]
    \centering
    \includegraphics[width=0.75\textwidth]{../surrogate-models/shap_decay_overlay.png}
    \caption{Overlay of normalised SHAP decay curves for all nine targets.
    Each curve is scaled so that its maximum equals 1.
    All targets show rapid importance decay, with most predictive signal concentrated in the top 20--30 features.}
    \label{fig:shap_decay_overlay}
\end{figure}


% =====================================================================
% FEATURE REPRESENTATION ABLATION
% =====================================================================

\section{Feature representation ablation}
\label{app:ablation}

To verify that raw RDKit descriptors represent the best choice of molecular representation for this task, we systematically compared three alternatives using the identical cluster-based splitting protocol and hyperparameter optimisation setup as the main models.

First, we evaluated Morgan circular fingerprints across nine configurations combining three radii ($r \in \{2, 3, 4\}$) and three bit-vector lengths ($b \in \{512, 1024, 2048\}$), each used as a raw binary fingerprint vector.

Second, we applied two PCA-based compression strategies to the same 217 RDKit descriptors.
In the \emph{global PCA} approach, all descriptors were standardised and projected onto principal components retaining 95\% of the variance, reducing the feature set to between 35 and 46 components depending on the target-specific molecule subset.

In the \emph{two-stage PCA} approach, descriptors were first partitioned into chemically interpretable groups based on naming conventions: topological indices (Chi, Kappa), five families of surface-area distribution bins (PEOE\_VSA, SMR\_VSA, SlogP\_VSA, EState\_VSA, VSA\_EState), molecular mass variants, ring count descriptors, BCUT2D global descriptors, and MQN fragment counts.
Within each group, a single principal component was extracted.
Groups for which this PC1 explained at least 80\% of within-group variance (in practice: Chi, Kappa, and MolWt families) were retained as single interpretable scalars;
all remaining groups and ungrouped descriptors were pooled and compressed by global PCA at 95\% variance.
This yielded 4 group-level features combined with 36--47 pool-level components per target, giving approximately 40--51 features in total.
Both PCA models were fitted independently per target property on the molecule subset used for that target.

Table~\ref{tab:ablation} summarises performance for three representative targets across all evaluated feature strategies.
Morgan circular fingerprints consistently underperform the RDKit-based approach across all targets and configurations.
The best-performing Morgan variant (radius 4, 1024 bits) achieves $R^2 = 0.246 \pm 0.142$ for pCMC, $R^2 = 0.506 \pm 0.055$ for $D_{\mathrm{SOL}}$, and $R^2 = 0.176 \pm 0.054$ for $\eta$ --- compared to $0.518 \pm 0.143$, $0.920 \pm 0.012$, and $0.818 \pm 0.044$ for the raw RDKit baseline.
The gap is most pronounced for transport properties, which are most sensitive to the explicit physicochemical features encoded by RDKit descriptors --- surface area distribution, partial charge distribution, molecular weight --- that circular fingerprints do not capture directly.

Both PCA-based variants of the RDKit descriptors yield performance slightly below the full 217-feature set.
The two-stage PCA ($R^2 = 0.518 \pm 0.115$ for pCMC, $0.845 \pm 0.044$ for $D_{\mathrm{SOL}}$, $0.772 \pm 0.084$ for $\eta$) narrows the gap relative to the global PCA approach ($0.471 \pm 0.149$, $0.825 \pm 0.039$, $0.730 \pm 0.072$) but neither recovers the full-descriptor performance.
The fact that compression leads to a measurable drop indicates that the 217 RDKit features do not introduce significant overfitting on this dataset, and that the raw descriptor set represents the best available balance of predictive performance and feature-level interpretability.

\begin{table}[ht]
\centering
\resizebox{\linewidth}{!}{%
\begin{tabular}{lcccccc}
\toprule
\textbf{Representation} & \multicolumn{2}{c}{\textbf{pCMC}} & \multicolumn{2}{c}{\textbf{$D_{\mathrm{SOL}}$}} & \multicolumn{2}{c}{\textbf{$\eta$}} \\
 & $R^2$ & $\rho$ & $R^2$ & $\rho$ & $R^2$ & $\rho$ \\
\midrule
Raw RDKit (217 features)                  & $0.518 \pm 0.143$ & $0.731 \pm 0.099$ & $0.920 \pm 0.012$ & $0.953 \pm 0.009$ & $0.818 \pm 0.044$ & $0.923 \pm 0.031$ \\
RDKit + two-stage PCA ($\sim$40--51 feat.)& $0.518 \pm 0.115$ & $0.726 \pm 0.098$ & $0.845 \pm 0.044$ & $0.912 \pm 0.023$ & $0.772 \pm 0.084$ & $0.911 \pm 0.022$ \\
RDKit + global PCA (35--46 PCs)           & $0.471 \pm 0.149$ & $0.702 \pm 0.115$ & $0.825 \pm 0.039$ & $0.901 \pm 0.020$ & $0.730 \pm 0.072$ & $0.887 \pm 0.046$ \\
Morgan ($r$=4, $b$=1024)                  & $0.246 \pm 0.142$ & $0.510 \pm 0.166$ & $0.506 \pm 0.055$ & $0.705 \pm 0.047$ & $0.176 \pm 0.054$ & $0.537 \pm 0.044$ \\
\bottomrule
\end{tabular}%
}
\caption{Feature representation ablation for three representative targets (pCMC: best experimental target;
$D_{\mathrm{SOL}}$: best MD target;
$\eta$: strong MD target).
All variants use the same cluster-based evaluation protocol and 30-trial HPO.
Only the best-performing Morgan configuration is shown;
results for all 9 Morgan variants are consistent with this trend.}
\label{tab:ablation}
\end{table}


% =====================================================================
% FEATURE-LEVEL ANALYSIS
% =====================================================================

\section{Feature-level analysis of structure--property relationships}
\label{app:feature_importance}

To investigate whether any of the RDKit features are particularly important for surfactant behaviour, permutation importance was used to identify the top-10 features in each of the 5 test sets for each of the 9 targets.
The importance was summed across the 5 test sets to obtain a top-10 list per target property.

A small number of descriptors recur across the nine targets: one feature appears in the top-10 list of six different models, and two further features appear in five each.
The three most recurrent descriptors are VSA\_EState9, which captures the distribution of molecular surface area across specific electronic environments;
fr\_NH0, reflecting the presence of tertiary amine functionalities;
and BCUT2D\_MRLOW, a global descriptor related to molecular polarisability and size.
Despite this recurrence, the overall overlap of important features across targets remains limited, indicating that no single descriptor universally governs all surfactant properties considered here;
rather, different properties depend on partially distinct aspects of molecular structure, consistent with the multifaceted nature of surfactant behaviour.
Notably, all three of the most-recurrent descriptors contribute to models predicting diffusion-related properties, suggesting that they capture structural characteristics relevant for molecular mobility in solution.

A complementary SHAP-based analysis is presented in Section~\ref{app:shap}.


% =====================================================================
% UNCERTAINTY ESTIMATION DIAGNOSTICS
% =====================================================================

\section{Uncertainty estimation diagnostics}
\label{app:unc_diagnostics}

Figure~\ref{fig:unc_nndist} shows absolute test error against nearest-neighbour (NN) distance to the training set in standardised RDKit feature space, for all nine target properties.
Spearman correlations are positive for most targets, with the strongest signal observed for pCMC ($\rho = 0.34$), pC$_{20}$ ($\rho = 0.34$), and $\eta$ ($\rho = 0.28$).
Correlations are weak or negligible for $\gamma_{\mathrm{CMC}}$ ($\rho = 0.04$) and $\gamma$ ($\rho = 0.11$), indicating that proximity in RDKit feature space is not predictive of prediction quality for these targets.
Notably, $D_{\mathrm{MOL}}$ shows a marginally negative correlation ($\rho = -0.06$), suggesting that the feature space does not capture the structural variation relevant for solute diffusion.

\begin{figure}[ht]
    \centering
    \includegraphics[width=\textwidth]{../surrogate-models/unc_error_vs_nndist.pdf}
    \caption{Absolute prediction error versus nearest-neighbour distance to the training set in standardised feature space, for all nine target properties.
    Each point corresponds to one test molecule;
    colours distinguish the five independent test splits (viridis scale).
    The Spearman rank correlation $\rho$ is annotated in each panel.
    Points with NN distance exceeding 40 are excluded;
    the number of excluded points is noted where non-zero.}
    \label{fig:unc_nndist}
\end{figure}

Figure~\ref{fig:unc_ensstd} shows absolute test error against ensemble standard deviation across all test splits.
Points below the identity line $y = x$ indicate an underestimate of error (overconfidence), while points above indicate an overestimate (underconfidence).
Spearman correlations are highest for $D_{\mathrm{MOL}}$ ($\rho = 0.45$), $A_{\min}$ ($\rho = 0.31$), pCMC ($\rho = 0.29$), $D_{\mathrm{SOL}}$ ($\rho = 0.28$), and $\eta$ ($\rho = 0.28$).
For $\gamma$ ($\rho = 0.06$) and $\gamma_{\mathrm{CMC}}$ ($\rho = 0.07$), ensemble disagreement carries almost no information about prediction error.
The complementarity of the two heuristics is illustrated by $D_{\mathrm{MOL}}$, for which NN distance fails ($\rho = -0.06$) but ensemble standard deviation is the most informative signal across all targets ($\rho = 0.45$).

\begin{figure}[ht]
    \centering
    \includegraphics[width=\textwidth]{../surrogate-models/unc_error_vs_ens_std.pdf}
    \caption{Absolute prediction error versus ensemble standard deviation for all nine target properties.
    Each point corresponds to one test molecule;
    colours distinguish the five independent test splits (viridis scale).
    The dashed diagonal line indicates perfect correspondence between ensemble standard deviation and absolute error.
    The Spearman rank correlation $\rho$ is annotated in each panel.}
    \label{fig:unc_ensstd}
\end{figure}


% =====================================================================
% CONFORMAL CALIBRATION ANALYSIS
% =====================================================================

\section{Conformal calibration analysis}
\label{app:conformal}

The conformal calibration curves presented in the main manuscript reveal a particularly informative pattern: the direction of the deviation from the diagonal is not the same for all properties, and this asymmetry is physically meaningful.

For the best-predicted MD-derived targets ($D_{\mathrm{SOL}}$, $\eta$) and for $\gamma_{\mathrm{CMC}}$, the calibration curves fall systematically below the diagonal, indicating under-coverage.
This occurs because the model has learned tight, accurate predictions in-distribution, resulting in small calibration residuals and narrow conformal intervals.
When these intervals are applied to test molecules from held-out structural clusters --- which are by construction more distant from the training distribution than the validation folds --- the intervals are insufficient to capture the inflated errors, and coverage falls below the nominal level.

By contrast, for $A_{\min}$, $\Gamma_{\max}$, and pC$_{20}$, the calibration curves sit consistently above the diagonal: the model is over-conservative, and requesting 90\% nominal coverage routinely yields 91--96\% empirical coverage.
These properties are difficult to predict \textit{uniformly} across chemical space rather than specifically at the boundary of the training distribution.
Because the calibration residuals already reflect this uniform difficulty, the resulting conformal intervals are wide enough to cover even the test-set molecules from novel clusters.
The over-coverage confirms that these properties are hard everywhere --- not harder in the held-out clusters than in the validation folds.

This has a practical implication for users of SurfPro-MD.
For MD-derived properties such as $\eta$ and $D_{\mathrm{SOL}}$, where under-coverage is most pronounced, a user requesting 90\% prediction intervals should expect approximately 85\% actual coverage on structurally novel molecules.
Requesting 95\% nominal coverage recovers approximately 92--93\% empirical coverage for these targets (see the uncertainty summary table in the main manuscript).
For $A_{\min}$ and $\Gamma_{\max}$, reported 90\% intervals are already conservative by several percentage points, and no correction is needed.
For pCMC and $D_{\mathrm{MOL}}$, where calibration curves lie essentially on the diagonal, the conformal intervals can be trusted at face value.

More broadly, the conformal calibration curve functions as a property-specific fingerprint of the degree to which cluster-based generalisation differs from interpolation.
Properties for which the test clusters are genuinely more OOD than the validation folds --- those with tight in-distribution fits and under-coverage --- are precisely the ones where the cluster-based splitting protocol is most stringent.
Properties for which the calibration curve exceeds the diagonal reflect targets where the whole chemical space is uniformly hard, and the splitting strategy does not introduce additional difficulty.
This distinction is not visible in aggregate performance metrics such as $R^2$ or Spearman $\rho$, making the calibration analysis a complementary diagnostic that reveals structure in the difficulty of the learning problem that aggregate metrics conceal.

\end{document}
